\begin{explanationbox}

\textbf{\textcolor{CExplanation}{一句话直觉}}

三角形面积公式 $2S=\sqrt{PQ+QR+RP}$ 将面积计算转化为三个辅助量 $P,Q,R$ 的代数运算，避免了对半周长和海伦公式中复杂根号的处理.

\medskip
\textbf{\textcolor{CExplanation}{核心拆解}}
\begin{description}[style=nextline, leftmargin=2.8em, labelsep=0.5em, itemsep=0.45em]
    \item[\textbf{要点一（对称构造）：}] 由 $P+Q=x^2$ 等对称方程组定义 $P,Q,R$，使得它们只依赖于边长平方，形式对称.
    \item[\textbf{要点二（代数关系）：}] $PQ+QR+RP = 4S^2$，因此 $2S = \sqrt{PQ+QR+RP}$.
\end{description}

\medskip
\textbf{\textcolor{CExplanation}{几何本质（海伦等价）}}

该公式与海伦公式完全等价，但将半周长 $s$ 的表达式替换为三边平方的交叉项，几何上对应 $P,Q,R$ 分别等于 $yz\cos A,\;zx\cos B,\;xy\cos C$。

\medskip
\textbf{\textcolor{CExplanation}{代数意义（平方对称）}}

公式中只出现边长平方和它们的乘积，不含边长本身的一次项，因此在处理根式边长时，可直接代入平方值，避免开方带来的嵌套根式.

\medskip
\textbf{\textcolor{CExplanation}{考点价值（计算简化）}}
\begin{itemize}[leftmargin=*, itemsep=0.5em, topsep=0.4em]
    \item \textbf{考法一（根式边长）：} 已知三边为 $\sqrt{a},\sqrt{b},\sqrt{c}$，直接令 $x^2=a$ 等，快速计算面积.
    \item \textbf{考法二（解三角形综合）：} 结合余弦定理、面积公式时，可转化为 $P,Q,R$ 的代数方程求解边长.
\end{itemize}

\medskip
\textbf{\textcolor{CExplanation}{顿悟点}}

注意到 $P = \frac{x^2+z^2-y^2}{2}$ 正是余弦定理中 $x^2+z^2-y^2 = 2xz\cos B$ 的 $\frac12$ 倍，因此 $P,Q,R$ 的几何本质是边长乘积与夹角的余弦组合.

\medskip
\textbf{\textcolor{CExplanation}{使用场景}}
\begin{itemize}[leftmargin=*, itemsep=0.5em, topsep=0.4em]
    \item \textbf{场景一（边长含根号）：} 三边长度包含 $\sqrt{\ }$ 时，用平方值代入避免二次根式嵌套.
    \item \textbf{场景二（已知平方关系）：} 已知三边平方和、两两平方差等关系时，直接构造 $P,Q,R$ 的方程组.
\end{itemize}

\end{explanationbox}
